\documentclass[11pt,a4paper]{article}
\usepackage{graphicx}
\usepackage{listings} % Required for beautiful code/terminal formatting
\usepackage{xcolor}   % Required for optional coloring in code blocks
\usepackage{microtype} % Improves text justification

% Optional: Configure the look of the terminal/service dumps
\lstset{
    basicstyle=\ttfamily\small,
    backgroundcolor=\color{gray!10},
    frame=single,
    breaklines=true,
    postbreak=\mbox{\textcolor{red}{$\hookrightarrow$}\space},
    keepspaces=true,
    showstringspaces=false
}

\begin{document}

% ==========================================
% 1. TITLE PAGE (First Page Only)
% ==========================================
\begin{titlepage}
    \centering
    \vspace*{2cm}
    
    {\large \textbf{Heilbronn University of Applied Sciences} \par}
    {\large Course: Real Time Systems \par}
    
    \vspace{2.5cm}
    
    {\huge \bfseries Open1722 CAN Latency Tests \par}
    \vspace{0.5cm}
    {\Large Project Report \par}
    
    \vfill
    
    {\large \textbf{Author:}} \\
    {\large Sripiranavan Yogarajah \par}
    
    \vfill
    
    {\large \today \par}
\end{titlepage}

% ==========================================
% 2. TABLE OF CONTENTS & PAGE NUMBERS SECTION
% ==========================================
\newpage
\pagenumbering{roman} % Use Roman numerals for intro pages if desired
\tableofcontents
\clearpage

% ==========================================
% 3. MAIN CONTENT
% ==========================================
\pagenumbering{arabic} % Switch to regular numbers (1, 2, 3...) for content

\section{Introduction}
This is the project report for the CAN Latency Tester. Below are the underlying environment configurations and systemd service specifications utilized during the testing setup.

\subsection{System Environment}
The test environment runs on the following OS distribution release details:

\begin{lstlisting}[caption={OS Release Information (/etc/os-release)}]
sripiranavan@sripiranavan:~ $ cat /etc/os-release 
PRETTY_NAME="Debian GNU/Linux 13 (trixie)"
NAME="Debian GNU/Linux"
VERSION_ID="13"
VERSION="13 (trixie)"
VERSION_CODENAME=trixie
DEBIAN_VERSION_FULL=13.4
ID=debian
HOME_URL="https://www.debian.org/"
SUPPORT_URL="https://www.debian.org/support"
BUG_REPORT_URL="https://bugs.debian.org/"
\end{lstlisting}

\subsection{Network and Bridge Configuration}
To initialize the Controller Area Network (CAN) interface at a bitrate of 250,000 bps, the following systemd service is deployed:

\begin{lstlisting}[caption={CAN1 Interface Service (can1.service)}]
sripiranavan@sripiranavan:/etc/systemd/system $ cat can1.service 
[Unit]
Description=Bring up CAN1 interface
After=network.target

[Service]
Type=oneshot
ExecStart=/sbin/ip link set can1 up type can bitrate 250000
ExecStop=/sbin/ip link set can1 down
RemainAfterExit=yes

[Install]
WantedBy=multi-user.target
\end{lstlisting}

\noindent Subsequently, the Open1722 CAN-to-Ethernet Bridge software is managed via the following service file to bind the `can1` interface to the network stack:

\begin{lstlisting}[caption={Open1722 Bridge Service (acf-can-bridge.service)}]
sripiranavan@sripiranavan:/etc/systemd/system $ cat acf-can-bridge.service 
[Unit]
Description=Open1722 CAN to Ethernet Bridge (Pi1)
After=can1.service network-online.target
Wants=can1.service network-online.target

[Service]
Type=simple
ExecStart=/sripiranavan/Open1722/build/examples/acf-can/linux/acf-can-bridge \
  --dst-addr e4:5f:01:2a:87:70 \
  -i eth0 \
  --canif can1

Restart=always
RestartSec=2

[Install]
WantedBy=multi-user.target
\end{lstlisting}


\section{Methodology}
This section outlines the experimental setup, detailing the specific hardware components and software frameworks utilized to construct the CAN Latency Tester environment.

\subsection{System Components and Hardware Setup}
The hardware architecture consists of multiple nodes configured to transmit, bridge, and monitor CAN frame latency. The complete bill of materials is detailed in Table~\ref{tab:hardware}.

\begin{table}[h]
    \centering
    \caption{Hardware Components}
    \label{tab:hardware}
    \begin{tabular}{|l|c|p{7cm}|}
        \hline
        \textbf{Component} & \textbf{Quantity} & \textbf{Description / Role in Project} \\ \hline
        Raspberry Pi 4 & 2 & Acts as the Linux host running the Open1722 CAN-to-Ethernet bridge application. \\ \hline
        STM32F439ZI & 3 & Microcontroller development boards used as dedicated CAN nodes for traffic generation or timestamping. \\ \hline
        CAN Transceiver & 4 & Interfaces the logic-level (TTL) signals from the controllers to the physical CAN bus differential voltages. \\ \hline
        120\,$\Omega$ Resistor & 4 & Used for proper bus termination networks to prevent signal reflections across the nodes. \\ \hline
        USB-CAN Adapter & 1 & Connected to a host machine for non-intrusive bus monitoring, packet sniffing, and verification. \\ \hline
    \end{tabular}
\end{table}

\subsection{Software Stack and Environment}
To ensure predictable execution, real-time networking, and standardized encapsulation, the specific software frameworks listed in Table~\ref{tab:software} were deployed.

\begin{table}[h]
    \centering
    \caption{Software Frameworks and Operating Systems}
    \label{tab:software}
    \begin{tabular}{|l|l|p{6.5cm}|}
        \hline
        \textbf{Software / OS} & \textbf{Version / Variant} & \textbf{Application Purpose} \\ \hline
        Debian GNU/Linux & 13 (Trixie) & Operating system for the Raspberry Pi 4 hosts, managing network routing and systemd background services. \\ \hline
        Zephyr OS & RTOS (Latest stable) & Real-Time Operating System running on the STM32 nodes to guarantee precise, deterministic CAN frame processing timelines. \\ \hline
        Open1722 & Custom Build & High-performance audio/video bridging transport protocol engine used to encapsulate CAN traffic over Ethernet layers. \\ \hline
    \end{tabular}
\end{table}


\subsection{Network Topology}
% Tip: You can describe how they are connected here. 
The nodes are physically linked over a shared bus terminated at both ends by the 120\,$\Omega$ resistors. The STM32 units running Zephyr OS interact natively with the CAN transceivers, while the Raspberry Pi 4 targets leverage the Open1722 infrastructure via systemd service units to bridge local data over the standard network interface setup.


\section{Results}
Present the results obtained from the CAN Latency Tester.

\begin{figure}[h]
    \centering
    % \includegraphics[width=0.8\textwidth]{path/to/your/image.png}
    \caption{A caption for the figure.}
    \label{fig:example}
\end{figure}


\section{Conclusion}
Summarize the findings and conclude the report.

\end{document}